\chapter{PHƯƠNG PHÁP LUẬN}

Trong chương này, chúng tôi sẽ trình bày và đề xuất phương pháp luận của nghiên cứu, bao gồm: pipeline tổng thể, quy trình tổng hợp dữ liệu sử dụng kết hợp hai chiến lược là few-shot và aspect-based Chain-of-Thought, hai giai đoạn huấn luyện là QLoRA theo tác vụ và tinh chỉnh toàn phần (Full-parameter fine-tuning), thư viện tăng tốc Unsloth và cuối cùng là các kỹ thuật tối ưu để toàn bộ quy trình kể trên chạy được trên môi trường notebook miễn phí như Kaggle hoặc Google Colab với runtime GPU T4 16GB.

\section{Pipeline tổng thể}

Phương pháp luận của chúng tôi được phát triển dựa trên kiến trúc hai giai đoạn LoRA-to-Full của đội MinLegal~\cite{luan2025minlegal} và ý tưởng phân rã theo khía cạnh (aspect-based) của Bosch@AI~\cite{quang2025bosch} vào pipeline sinh dữ liệu. Phương pháp luận được tổ chức và phân chia thành năm giai đoạn liên tiếp như minh họa ở Hình~\ref{fig:pipeline}. Giai đoạn đầu tiên là chuẩn bị dữ liệu gốc: đầu tiên chúng tôi tải 440 mẫu từ ba subset, tiếp đó chúng tôi thực hiện phân tích khám phá (EDA) và chuyển đổi sang định dạng hội thoại của Qwen. Giai đoạn thứ hai là tổng hợp dữ liệu, như nhắc trên kết hợp hai chiến lược là few-shot (A) và aspect-based (B) để xây dựng một tập dữ liệu huấn luyện gồm khoảng 2,346 mẫu sinh. Giai đoạn thứ ba là chúng tôi bắt đầu huấn luyện ba adapter QLoRA riêng biệt ứng với lần lượt ba tác vụ theo đề ra của thử thách. Giai đoạn thứ tư, sau khi huấn luyện chúng tôi hợp nhất ba adapter bằng nội suy tuyến tính (Linear Interpolation) rồi bắt đầu tinh chỉnh toàn phần mô hình đã hợp nhất. Giai đoạn cuối cùng của phương pháp luận là chúng tôi sẽ thực nghiệm và đánh giá mô hình trên ba tác vụ của thử thách như kể trên.

\begin{figure}[ht]
\centering
\begin{tikzpicture}[
    font=\footnotesize,
    box/.style={rectangle, draw, rounded corners, align=center, minimum height=1cm, text width=2.3cm, fill=blue!5},
    arr/.style={->, thick}
]
\node[box] (d) at (0,0) {1. Chuẩn bị\\dữ liệu gốc\\(440 mẫu)};
\node[box, fill=orange!10] (g) at (3,0) {2. Sinh dữ liệu\\A + B\\($\sim$2.346 mẫu)};
\node[box, fill=green!8] (s1) at (6,0) {3. Stage 1\\QLoRA\\3 adapter};
\node[box, fill=green!8] (m) at (9,0) {4. Hợp nhất +\\Stage 2\\Full FT};
\node[box, fill=violet!8] (e) at (12,0) {5. Đánh giá\\3 tác vụ};
\draw[arr] (d) -- (g);
\draw[arr] (g) -- (s1);
\draw[arr] (s1) -- (m);
\draw[arr] (m) -- (e);
\end{tikzpicture}
\caption{Pipeline tổng thể năm giai đoạn của phương pháp đề xuất.}
\label{fig:pipeline}
\end{figure}

\section{Giai đoạn sinh dữ liệu kết hợp}

Một trong số các điểm mới của nghiên cứu nằm ở việc kết hợp hai chiến lược sinh dữ liệu vốn ban đầu được áp dụng riêng rẽ trong hai phương pháp gốc, bây giờ chúng tôi áp dụng trên cùng một baseline và cùng điều kiện huấn luyện. Một điểm mới nữa, khác với MinLegal (sử dụng Gemini 2.5 Flash) và Bosch@AI (sử dụng GPT-4o), báo cáo này sử dụng \textbf{Anthropic Claude Haiku 4.5}~\cite{anthropic2025claude} làm mô hình giáo viên (teacher model) cho cả hai chiến lược vì mô hình này có chất lượng sinh văn bản tiếng Việt tốt và khả năng tuân thủ định dạng đầu ra ổn định. Hình~\ref{fig:datagen} mô tả tổng quan quy trình sinh dữ liệu theo hai nhánh.

\begin{figure}[ht]
\centering
\begin{tikzpicture}[
    font=\footnotesize,
    box/.style={rectangle, draw, rounded corners, align=center, minimum height=0.9cm, text width=3.2cm},
    io/.style={box, fill=gray!10},
    proc/.style={box, fill=blue!6},
    arr/.style={->, thick}
]
% Input
\node[io] (seed) at (0,0) {440 mẫu gốc + văn bản luật};
% Strategy A branch (top)
\node[proc, fill=orange!10] (a1) at (4.5,1.3) {Chiến lược A: chọn 5 mẫu few-shot};
\node[proc, fill=orange!10] (a2) at (9,1.3) {Sinh MC/NLI mới (Claude Haiku 4.5)};
% Strategy B branch (bottom)
\node[proc, fill=teal!10] (b1) at (4.5,-1.3) {Chiến lược B: trích khía cạnh (aspect)};
\node[proc, fill=teal!10] (b2) at (9,-1.3) {Sinh CoT tam đoạn luận theo aspect};
% Merge
\node[io, fill=violet!8, text width=2.6cm] (merge) at (12.8,0) {Khử trùng lặp + gộp ($\sim$2.346)};
\draw[arr] (seed) -- (a1);
\draw[arr] (seed) -- (b1);
\draw[arr] (a1) -- (a2);
\draw[arr] (b1) -- (b2);
\draw[arr] (a2) -- (merge);
\draw[arr] (b2) -- (merge);
\end{tikzpicture}
\caption{Quy trình sinh dữ liệu kết hợp hai chiến lược. Nhánh trên (A) sinh dữ liệu MC/NLI bằng few-shot, nhánh dưới (B) sinh dữ liệu tam đoạn luận (Syllogism) bằng aspect-based Chain-of-Thought.}
\label{fig:datagen}
\end{figure}

\subsection{Chiến lược A - Few-shot (theo MinLegal)}

Chiến lược A dựa trên prompting few-shot: Đầu tiên, với mỗi tác vụ, hệ thống sẽ chọn ngẫu nhiên năm mẫu gốc làm ví dụ minh họa. Tiếp đó, hệ thống sẽ đưa prompt vào cùng với yêu cầu mô hìnhInput
\node[io] (seed) at (0,0) {440 mẫu gốc + văn bản luật};
% Strategy A branch (top)
\node[proc, fill=orange!10] (a1) at (4.5,1.3) {Chiến lược A: chọn 5 mẫu few-shot};
\node[proc, fill=orange!10] (a2) at (9,1.3) {Sinh MC/NLI mới (Claude Haiku 4.5)};
% Strategy B branch (bottom)
\node[proc, fill=teal!10] (b1) at (4.5,-1.3) {Chiến lược B: trích khía cạnh (aspect)};
\node[proc, fill=teal!10] (b2) at (9,-1.3) {Sinh CoT tam đoạn luận theo aspect};
% Merge
\node[io, fill=violet!8, text width=2.6cm] (merge) at (12.8,0) {Khử trùng lặp + gộp ($\sim$2.346)};
\draw[arr] (seed) -- (a1);
\draw[arr] (seed) -- (b1);
\draw[arr] (a1) -- (a2);
\draw[arr] (b1) -- (b2);
\draw[arr] (a2) -- (merge);
\draw[arr] (b2) -- (merge);
\end{tikzpicture}
\caption{Quy trình sinh dữ liệu kết hợp hai chiến lược. Nhánh trên (A) sinh dữ liệu MC/NLI bằng few-shot, nhánh dưới (B) sinh dữ liệu tam đoạn luận (Syllogism) bằng aspect-based Chain-of-Thought.}
\label{fig:datagen}
\end{figure}

\subsection{Chiến lược A - Few-shot (theo MinLegal)}

Chiến lược A dựa trên prompting few-shot: Đầu tiên, với mỗi tác vụ, hệ thống sẽ chọn ngẫu nhiên năm mẫu gốc làm ví dụ minh họa. Tiếp đó, hệ thống sẽ đưa prompt vào cùng với yêu cầu mô hình sinh thêm các câu hỏi-đáp án mới tương tự về cấu trúc nhưng về chủ đề pháp lý khác nhau (ví dụ như thuế, đất đai, lao động, giao thông, hình sự, dân sự...). Tiếp theo, các kết quả được sinh ra sẽ được ban đầu lọc tự động (ví dụ như loại bỏ mẫu sai cấu trúc) và kết hợp với rà soát thủ công. Chúng tôi thấy chiến lược này tốt về việc mở rộng độ phủ (bề rộng) của dữ liệu, do đó có thể thấy chiến lược đã kế thừa trực tiếp được cách làm mà đã giúp đội MinLegal mở rộng từ 440 lên 3,537 mẫu dữ liệu.

\subsection{Chiến lược B - Aspect-based Chain-of-Thought (theo Bosch)}

Chiến lược B mô phỏng quy trình hai bước của Bosch~\cite{quang2025bosch}. \textbf{Bước B.1 (trích khía cạnh (aspect))}: Với mỗi điều luật, mô hình giáo viên xác định 1-3 khía cạnh pháp lý trọng yếu trong khối văn bản (corpus) (ví dụ điều kiện áp dụng, mức xử phạt, phạm vi điều chỉnh, hay ngoại lệ,..). \textbf{Bước B.2 (sinh CoT)}: Với mỗi khía cạnh vừa mới sinh, mô hình sinh một tình huống pháp lý kèm chuỗi suy luận (Chain-of-Thought) đặt trong thẻ tư duy \texttt{<think>...</think>}, tiếp đó sinh cấu trúc tam đoạn luận theo thử thách bao gồm tiền đề lớn, tiền đề nhỏ và kết luận. Chúng tôi thấy cách phân rã corpus theo các khía cạnh đã dạy mô hình phân biệt được các ngữ cảnh pháp lý một cách sâu sắc hơn (ví dụ: ranh giới giữa xử phạt hành chính với trách nhiệm hình sự), do đó nâng cao chất lượng lập luận cho tác vụ tam đoạn luận.

\subsection{Cấu trúc prompt}

Dưới đây là các prompt được sử dụng trong quy trình sinh dữ liệu và đánh giá:

\subsubsection{Prompt chiến lược A (Few-shot)}

Prompt few-shot cho tác vụ trắc nghiệm và NLI có cấu trúc như sau:

\begin{verbatim}
Bạn là chuyên gia tạo dữ liệu pháp luật Việt Nam. Dưới đây là
các ví dụ về [mô tả tác vụ].

VÍ DỤ 1:
ĐỀ: [nội dung câu hỏi mẫu 1]
ĐÁP: [đáp án mẫu 1]
...
VÍ DỤ 5:
ĐỀ: [nội dung câu hỏi mẫu 5]
ĐÁP: [đáp án mẫu 5]

Hãy tạo 6 mẫu MỚI cùng định dạng, ĐA DẠNG lĩnh vực pháp luật
(thuế, đất đai, lao động, giao thông, bảo hiểm xã hội, hình sự,
dân sự, hôn nhân gia đình, doanh nghiệp...).

Yêu cầu:
- Nội dung phải chính xác theo pháp luật Việt Nam hiện hành
- Đa dạng về độ khó (dễ, trung bình, khó)
- Không trùng lặp với các ví dụ đã cho

Trả về JSON array, mỗi phần tử có 2 key:
- "user": nội dung đề bài
- "assistant": đáp án
\end{verbatim}

\subsubsection{Prompt chiến lược B (Aspect-based CoT)}

\textbf{Prompt trích khía cạnh (B.1)}:
\begin{verbatim}
Dựa trên văn bản pháp luật được cung cấp, hãy xác định 1-3
khía cạnh pháp lý riêng biệt, điểm chính hoặc chủ đề được
bao quát (ví dụ: điều kiện áp dụng, mức xử phạt, phạm vi
điều chỉnh, ngoại lệ).

Nội dung văn bản:
"""[văn bản pháp luật]"""

Trả về CHỈ MỘT đối tượng JSON: { "aspects": ["Khía cạnh 1",
"Khía cạnh 2", ...] }
\end{verbatim}

\textbf{Prompt sinh suy luận (B.2)}:
\begin{verbatim}
Bạn là chuyên gia pháp luật Việt Nam chuyên tổng hợp dữ liệu.
Dựa trên văn bản pháp luật và các khía cạnh đã xác định, hãy
tạo các ví dụ tam đoạn luận pháp lý.

Văn bản pháp luật:
"""[văn bản pháp luật]"""

Các khía cạnh pháp lý: [danh sách khía cạnh]

Với MỖI khía cạnh, tạo một câu hỏi tình huống pháp lý mở và
lời giải có cấu trúc:
- Phần phân tích chi tiết đặt trong thẻ <think>...</think>
- Sau đó lần lượt ghi:
  Tiền đề lớn: <quy phạm pháp luật chung>
  Tiền đề nhỏ: <tình huống cụ thể>
  Kết luận: <hệ quả pháp lý có căn cứ>

Trả về JSON array, mỗi phần tử có 2 key:
- "question": tình huống pháp lý cần phân tích
- "reasoning": lời giải đầy đủ
\end{verbatim}

\subsubsection{Prompt đánh giá tam đoạn luận (LLM-as-Judge)}

Prompt dùng mô hình Claude Haiku 4.5 làm giám khảo cho nhiệm vụ tam đoạn luận:
\begin{verbatim}
Bạn là giám khảo chấm điểm câu trả lời tam đoạn luận pháp lý.
Cho điểm từ 0.0 đến 1.0 dựa trên tính đúng đắn của tiền đề
lớn, tiền đề nhỏ và tính hợp lý của kết luận.
Chỉ trả về một con số, không giải thích.

Câu hỏi:
[tình huống pháp lý]

Câu trả lời:
[câu trả lời của mô hình]
\end{verbatim}

Việc bắt ràng buộc đầu ra định dạng JSON giúp tự động hóa khâu trích xuất và lọc tiếp theo, giảm thiểu can thiệp thủ công, tăng hiệu suất.

\subsection{Hợp nhất và kiểm soát chất lượng}

Dữ liệu sinh từ hai chiến lược được gộp lại, tiếp đó khử trùng lặp ngữ nghĩa, cân bằng số lượng theo tác vụ. Sau đó dữ liệu được chuẩn hóa về định dạng hội thoại của Qwen với một system prompt thống nhất. Để bảo đảm tính chính xác của việc đánh giá dữ liệu, tập dữ liệu chúng tôi sinh được qua mô hình trên sẽ được đánh giá tách biệt hoàn toàn khỏi 440 mẫu public test nhằm tránh rò rỉ dữ liệu.

\subsection{Cấu hình kỹ thuật sinh dữ liệu}

Chúng tôi chạy riêng chương trình sinh dữ liệu độc lập, không ảnh hưởng đến chương trình rèn luyện mô hình, và chương trình này được chạy trước khi huấn luyện, với mục đích tạo ra các file JSON chứa các dữ liệu tổng hợp. Cấu hình kỹ thuật như sau:

\begin{itemize}
    \item \textbf{Mô hình sinh}: Claude Haiku 4.5 (\texttt{claude-haiku-4-5}) qua Anthropic API, được chọn vì chất lượng sinh văn bản tiếng Việt tốt và khả năng tuân thủ định dạng JSON ổn định.
    \item \textbf{Giới hạn tốc độ}: $\sim$1 giây giữa các lần gọi API kèm cơ chế retry với backoff lũy thừa.
    \item \textbf{Số mẫu sinh}: Sau khi lọc các mẫu không hợp lệ về đáp án/cấu trúc và khử trùng lặp, thu được số lượng các mẫu như sau: MC - 894, NLI - 900, Syllogism - 552 (giới hạn bởi số văn bản nguồn của chiến lược B), tổng cộng 2,346 mẫu sinh. Toàn bộ 2,346 mẫu sinh tạo thành tập huấn luyện; 440 mẫu gốc được giữ riêng làm tập đánh giá.
    \item \textbf{Cơ chế checkpoint}: Lưu trạng thái mới nhất sau mỗi batch để tiếp tục được nếu bị gián đoạn (ví dụ: phiên notebook timeout, lỗi API,...).
    \item \textbf{Khử trùng lặp}: Loại bỏ mẫu trùng dựa trên băm nội dung đã chuẩn hóa.
\end{itemize}

Việc tách riêng sinh dữ liệu khỏi huấn luyện theo chúng tôi mang lại hai lợi ích chính. Thứ nhất, có thể kiểm tra, rà soát và chỉnh sửa thủ công dữ liệu tổng hợp trước khi huấn luyện. Thứ hai, giúp tiết kiệm thời gian khi cần chạy lại huấn luyện mà không phải sinh lại dữ liệu.

\section{Giai đoạn 1 - QLoRA theo từng tác vụ}

Ở giai đoạn này, chúng tôi huấn luyện ba adapter QLoRA độc lập, mỗi adapter chuyên cho một trong ba tác vụ của thử thách. Mỗi adapter được áp dụng lên bảy ma trận chiếu của khối Transformer (\texttt{q\_proj}, \texttt{k\_proj}, \texttt{v\_proj}, \texttt{o\_proj}, \texttt{gate\_proj}, \texttt{up\_proj}, \texttt{down\_proj}) với cấu hình hạng $r=16$, hệ số tỉ lệ $\alpha=32$ và tỉ lệ dropout là 0.1. Tiếp theo, Mô hình nền được lượng tử hóa 4-bit NF4 thông qua thư viện Unsloth, theo đúng cơ chế QLoRA~\cite{dettmers2023qlora}. Tốc độ học ở giai đoạn này là $2 \times 10^{-4}$, huấn luyện trong 3 epoch với độ dài chuỗi tối đa 2048 token. Chúng tôi tách riêng từng adapter nhằm giúp mỗi tác vụ học được đặc trưng chuyên biệt của mỗi tác vụ mà không bị nhiễu chéo từ các tác vụ khác.

Sau khi chọn checkpoint tốt nhất cho từng tác vụ, ba adapter được hợp nhất bằng \emph{nội suy tuyến tính} như sau: trọng số của adapter gộp là trung bình có trọng số của ba adapter thành phần. Do mỗi adapter chỉ tham số hóa lượng cập nhật hạng thấp $\Delta W = BA$~\cite{hu2022lora}, việc nội suy tuyến tính trên không gian tham số adapter là khả thi và rẻ về tính toán.

\section{Giai đoạn 2 - Tinh chỉnh adapter hợp nhất}

Giai đoạn 2 bắt đàu từ adapter đã hợp nhất ở giai đoạn 1 và tiếp tục huấn luyện adapter này trên toàn bộ tập dữ liệu gộp trong 2 epoch nhờ điểm khởi tạo tốt từ giai đoạn 1, với tốc độ học $5 \times 10^{-5}$ - thấp hơn giai đoạn 1 để tránh phá vỡ tri thức đã học. Mục đích của giai đoạn này là "hòa giải" ba năng lực riêng lẻ thành một mô hình thống nhất. Một điểm khác so với cấu trúc của đội MinLegal~\cite{luan2025minlegal} là do giới hạn bộ nhớ GPU T4 16GB, chúng tôi không thực hiện huấn luyện tinh chỉnh toàn bộ tham số như gốc mà vẫn duy trì cơ chế LoRA ở cả hai giai đoạn.
